\section{A cumulant bound for the Fisher deficit}\label{sec:fisher}

We first estimate how far the Fisher information of a semicircularly
perturbed tuple is from that of a semicircular family with the same
covariance. The estimate will use the squared sizes of the higher
cumulant arrays, which will allow us to exploit the sparse cumulants
of the signal in the next section.

For a centered bounded self-adjoint $n$-tuple $C$, write $\kappa_m(C)$
for its order-$m$ free-cumulant array, with
\[
 \|\kappa_m(C)\|_2^2
 =\sum_{i_1,\ldots,i_m=1}^n
   |\kappa_m(C_{i_1},\ldots,C_{i_m})|^2.
\]
We use the moment--cumulant formula and its M\"obius inversion on
noncrossing partitions, the vanishing of mixed cumulants between free
algebras, and the fact that a standard semicircular family has only
second cumulants, given by its covariance
\cite[Proposition~11.4, Theorem~11.16, and Examples~11.21]{NS}.

The following testing identity converts the higher-cumulant terms into
an orthogonal series. For a bounded self-adjoint tuple $Z$, write
$E[\,\cdot\mid W^*(Z)]$ for the trace-preserving conditional
expectation, also viewed as the orthogonal
projection onto $L^2(W^*(Z),\tau)$ when its argument belongs only to $L^2$.

\begin{lemma}[Wick testing]\label{lem:wick}
Let $C$ be a centered bounded self-adjoint $n$-tuple, let $S$ be a
standard semicircular family free from $C$, and put $Z=C+\sqrt u S$,
where $u>0$. For a multi-index $\mathbf j=(j_1,\ldots,j_r)$, $r\geq1$,
and a monomial $P$ in the indeterminates $x_1,\ldots,x_n$, define
\[
 D_{\mathbf j}(P)
 =\sum_{P=P_0x_{j_1}P_1\cdots x_{j_r}P_r}
       \prod_{a=0}^r\tau(P_a(Z)).
\]
The sum selects increasing positions with the prescribed indices;
empty gap words $P_a$ contribute one. Extend $D_{\mathbf j}$ linearly
to polynomials. There is an orthonormal family of semicircular Wick
polynomials $U_{\mathbf j}(S)$, indexed by all such multi-indices, with
\begin{equation}\label{eq:wicktest}
 D_{\mathbf j}(P)=u^{-r/2}\tau(U_{\mathbf j}(S)^*P(Z)).
\end{equation}
In particular, the conjugate variables of $Z$ exist and are
\[
 \xi_i=u^{-1/2}E[S_i\mid W^*(Z)].
\]
\end{lemma}

\begin{proof}
Define $U_{\mathbf j}$ from the word $S_{j_1}\cdots S_{j_r}$ by summing
over sets of disjoint adjacent pairs in the \emph{original} index list:
delete each chosen pair with factor $-\delta_{j_a,j_{a+1}}$.
The empty set gives the original word. This definition also gives
$U_{\mathbf j}^*=U_{(j_r,\ldots,j_1)}$.

Expand $\tau(U_{\mathbf j}^*P(Z))$ by noncrossing cumulants. After
multilinearly expanding $Z=C+\sqrt u S$, freeness and semicircularity
show that every block meeting an initial $S$ position is a pair.
Its covariance is $\delta_{ij}$ against another initial $S$ and
$\sqrt u\,\delta_{ij}$ against a selected $S$ term from $Z$.

To account for the deletion terms, fix a choice of $C$ or $S$ at every
$P$-position in the multilinear expansion, and a noncrossing partition
$\pi$ on the original positions. Let $A(\pi)$ be its initial-initial pair
blocks whose endpoints are originally adjacent. For these fixed choices,
contributions correspond exactly to subsets $F\subseteq A(\pi)$: delete the pairs
in $F$ and restore them after expanding the remaining word. Restoring
adjacent pairs preserves noncrossingness and the covariance factors,
so their total coefficient is
\[
 \sum_{F\subseteq A(\pi)}(-1)^{|F|}=(1-1)^{|A(\pi)|}.
\]
If $\pi$ has any initial-initial pair, an innermost one is originally
adjacent. Indeed, its intervening positions are all initial, and each
must pair inside that interval to avoid crossing. Thus $A(\pi)$ is
nonempty and the partition cancels. This argument includes nested
internal pairs and does not require deleting newly adjacent positions.

In each surviving partition all initial positions pair externally with
positions of $P$, in reverse order. The initial indices in
$U_{\mathbf j}^*$ are $j_r,\ldots,j_1$, so the selected indices in $P$
occur as $j_1,\ldots,j_r$. The external pair arches prevent any remaining
block from connecting distinct gaps. Summing the unrestricted partitions
and the choices of $C$ or $S$ within each gap therefore gives its trace,
even though the entries of $C$ need not be free. Each external pair contributes $\sqrt u$, proving
\eqref{eq:wicktest}.

Apply this calculation with zero signal and $u=1$. A Wick word tests
to zero against a shorter word and gives the corresponding Kronecker
delta against a word of its own degree. Since every other term of a
Wick polynomial has smaller degree, this proves orthonormality against
Wick words of equal or smaller degree; conjugate symmetry covers the
remaining cases. For $r=1$, $D_i(P)=(\tau\otimes\tau)(\partial_iP(Z))$,
so \eqref{eq:wicktest} and conditional expectation give the stated
conjugate variables.
\end{proof}

\begin{proposition}[Squared-cumulant Fisher bound]\label{prop:fisher}
Let $C$ be a centered bounded self-adjoint $n$-tuple with diagonal
covariance $\tau(C_iC_j)=\delta_{ij}c_i$, where $c_i\geq0$, and let $S$ be a standard
semicircular family free from $C$. Put $V_u=C+\sqrt u S$ and
$v_i(u)=u+c_i$ for $u>0$. For every $u>0$,
\[
 \sum_i\frac1{v_i(u)}\leq\Phi^*(V_u)\leq\frac n u.
\]
Whenever
$\sum_{r\geq2}u^{-r}\|\kappa_{r+1}(C)\|_2^2<\infty$,
\begin{equation}\label{eq:fisherbound}
 0\leq\Phi^*(V_u)-\sum_i\frac1{v_i(u)}
 \leq\frac1{u^2}\sum_{r\geq2}u^{-r}\|\kappa_{r+1}(C)\|_2^2.
\end{equation}
\end{proposition}

\begin{proof}
Let $Z=V_u$ and let $\xi_i$ be supplied by Lemma~\ref{lem:wick}.
Contraction of conditional expectation gives $\Phi^*(Z)\leq n/u$.
Also $\tau(Z_i\xi_i)=1$ by testing the conjugate identity with $Z_i$.
Since $\|Z_i\|_2^2=v_i(u)$, Cauchy--Schwarz gives
$\|\xi_i\|_2^2\geq v_i(u)^{-1}$ for every $u>0$.

For the upper deficit bound, expand the block containing the first
position in $\tau(Z_iP(Z))$. Its other positions select the letters
in $D_{\mathbf j}(P)$, while its gaps contribute their moments. Thus
\[
 \tau(Z_iP(Z))=\sum_{r\geq1,\mathbf j}
 \kappa_{r+1}(Z_i,Z_{j_1},\ldots,Z_{j_r})D_{\mathbf j}(P).
\]
Only $r\leq\deg P$ contributes. Centering removes the order-one
cumulant, and diagonal covariance makes the order-two term
$v_i(u)D_i(P)$. Adding free semicircular noise leaves all higher
cumulants equal to those of $C$.

Define the ambient $L^2$ series
\[
 R_i=\sum_{r\geq2,\mathbf j}u^{-r/2}
 \kappa_{r+1}(C_i,C_{j_1},\ldots,C_{j_r})U_{\mathbf j}(S)^*.
\]
The adjoint Wick family is orthonormal, and the assumed summability
gives convergence in $L^2$ with
\[
 \sum_i\|R_i\|_2^2=\sum_{r\geq2}u^{-r}\|\kappa_{r+1}(C)\|_2^2.
\]
For a fixed polynomial $P$, all terms with $r>\deg P$ have zero trace
pairing by \eqref{eq:wicktest}. Moreover, $A\mapsto\tau(AP(Z))$ is
continuous on $L^2$. The finite cumulant expansion above therefore gives
\[
 \tau\big((Z_i-v_i(u)\xi_i)P(Z)\big)=\tau(R_iP(Z)).
\]
The coefficients in $R_i$ need no complex conjugation: this identity
uses the bilinear trace pairing, rather than a Hilbert inner product.
Testing with adjoint polynomials converts it into Hilbert-space
orthogonality. Polynomial $L^2$ density then identifies
\[
 Z_i-v_i(u)\xi_i=E[R_i\mid W^*(Z)].
\]

The conjugate variables are self-adjoint by their conditional-expectation
formula. Expanding the residual norm using $\tau(Z_i\xi_i)=1$ yields
\[
 \|Z_i-v_i(u)\xi_i\|_2^2
 =v_i(u)^2\big(\|\xi_i\|_2^2-v_i(u)^{-1}\big).
\]
Summing after division by $v_i(u)^2$, using projection contraction and
$v_i(u)\geq u$, proves \eqref{eq:fisherbound}.
\end{proof}

The proposition reduces the large-noise entropy estimate to control of
the cumulant arrays of the signal. Its separate bound $\Phi^*(V_u)\leq n/u$
will also ensure integrability on every compact interval of positive
noise times.
